\section{Conclusion}
\label{sec:conclusion}

% Wang T10 shape (CLAIM -> METHODOLOGY -> EVIDENCE -> CONSEQUENCE), one tight
% paragraph. Numbers mirror results.tex Summary (sec:results:summary) and the
% abstract; keep in sync.
This paper presents \approach, the first training-free, multi-stage \ac{LLM} workflow for architecture document to model traceability link recovery.
\approach derives an alias table once, then reuses it in two linkers that reach three reference
forms: a whole name, one word of a name, and an implicit reference.
Each linker sits behind an \ac{LLM}-based judge that rejects links with no supporting evidence, and
the two written forms are judged by one rule that reads the evidence their match computed.
Alongside it we introduce an architecture-driven evaluation suite that adds a per-component tail to the standard link-level metric on doc-code, and the \cmrname{} (\cmr) on doc-model.
On the established ARDoCo benchmark, \approach{} reaches a doc-model \avgfone of $0.952$, $14.2$pp above the strongest baseline run on the same backend, and $0.965$ \avgftwo, $13.7$pp above it,
 and carries this gain through to a doc-code \avgfone of $0.918$ ($+10.6$pp).
The standard metric is where this lead is smallest.
Under the size-aware suite it widens to $+35.3$pp worst-component \fone\ and $+30.6$pp harmonic-mean per-component \fone.
\approach{} worst component still sits at \fone\ $0.82$, well below the link-level \fone, even though it no longer abandons any documented component outright.
So the suite earns its place by keeping this headroom visible, where the standard metric would call the task solved.